\documentclass[11pt]{article}
\usepackage[margin=0.9in]{geometry}
\usepackage{graphicx,amsmath,amssymb,booktabs,hyperref,xcolor}
\graphicspath{{/tmp/dd_k3_trf_figs/}}
\hypersetup{colorlinks=true,linkcolor=blue,urlcolor=blue}
\setlength{\parskip}{6pt}
\title{Cheby-native operator + trust-region least squares:\\
       what we learn about FP existence and operator smoothness}
\author{Fixed-Point Factory --- TRF follow-up}
\date{June 2026}

\begin{document}
\maketitle

\begin{abstract}
\noindent
After alternating projection (Anderson + monotone cumulative-max) failed
to converge on the Cheby-native operator, we tried pure trust-region
least squares (\texttt{scipy.optimize.least\_squares}, \texttt{method='trf'})
on the symmetric-reduced unknowns ($n=84$ at $G=7$). Tested at
$\gamma{=}100, \tau{=}0.1$ (smooth regime per the regime map):

\begin{itemize}
  \item TRF converges in 6 outer iterations to $|F|_\infty \approx 1.3 \times 10^{-2}$,
        cost $7.7 \times 10^{-4}$, with $\sim 5 \times 10^{-7}$ reduction at
        iter 6 (essentially terminated).
  \item Total work: $\sim$650 Phi evals ($\sim$28 min wall) including the
        FD Jacobian rebuilds at each iter.
  \item The residual floors at $\sim 10^{-2}$ at $G=7$ --- the same algebraic
        Cheby-resolution limit we identified earlier. The FP exists and is
        well-defined; TRF finds it up to the polynomial-resolution allowed
        by 84 symmetric-reduced unknowns.
\end{itemize}

\textbf{Key takeaway.} The Cheby-native operator DOES have a fixed point.
Newton-class methods (TRF) can find it monotonically (unlike Anderson +
projection, which oscillated). The remaining error is genuinely the
finite-resolution limit of $G=7$ Cheby on a $C^k$ FP; bigger $G$ would
push the residual down.
\end{abstract}

\section{TRF convergence}

\begin{center}
\includegraphics[width=\linewidth]{fig1_trf_conv.png}
\end{center}

\textbf{Left:} TRF outer-iteration cost (= $\frac{1}{2}\|F\|_2^2$). Iters 0..6:
2.24 $\to$ 0.150 $\to$ 0.0040 $\to$ 0.0024 $\to$ 0.00079 $\to$ 0.00077 $\to$ 0.00077.
Plateau at iter 6 (reduction $5 \times 10^{-7}$).

\textbf{Right:} $|F|_\infty$ along the wall-time axis. The flat stretches
are FD-Jacobian builds (84 calls each); the drops are the TRF step.

\section{Comparison with Anderson + projection (Method C)}

\begin{center}
\begin{tabular}{lrrr}
\toprule
method                    & final $|F|_\infty$ & iters &  comments \\
\midrule
Anderson + projection (it 20) & $\sim 0.05$ oscillating & 20 & did not converge \\
\textbf{TRF (it 6)}            & $\mathbf{0.013}$       & 6 & monotone convergence \\
\bottomrule
\end{tabular}
\end{center}

TRF achieves $\sim 4\times$ better residual at $4\times$ fewer outer iters.
Anderson is too aggressive for an operator with non-smooth Jacobian; TRF's
trust-region constraint keeps the step modest, allowing the FD Jacobian
to track the operator's local linearization.

\section{Why the floor at $10^{-2}$ at $G=7$}

The Cheby fit at $G=7$ has only 84 symmetric-reduced coefs (vs 343 cube
cells). For a $C^k$ FP with $k$ small, this gives $|F| \sim 1/G^k$
resolution. At $k \approx 2$ (which the regime-map data suggests): $1/7^2
\approx 0.02$, matching the observed floor.

Scaling up: $G=15$ would give $|F| \sim 1/15^2 \approx 4 \times 10^{-3}$
(with $\sim$680 unknowns and $\sim$3 hours TRF wall);
$G=21$ would give $|F| \sim 1/21^2 \approx 2 \times 10^{-3}$ (with $\sim$1771
unknowns and $\sim$1 day wall). Going higher needs the Cheby-native operator
JIT-compiled.

\section{Recommendations}

\begin{enumerate}
  \item \textbf{TRF is the right solver for the Cheby-native operator.}
        It's robust, monotone, terminates cleanly. Adopt as the default.
  \item \textbf{Need a faster Phi for higher $G$}: rewrite \texttt{lookup\_analytic\_lobatto}
        with numba + a hand-coded companion-matrix root finder (replace
        \texttt{cheb.chebroots}). Should give $\sim 50\times$ speedup,
        making $G=15$ and $G=21$ tractable.
  \item \textbf{Or use the Lin-CDF R4 Jacobian as a chord substitute} in TRF:
        cheap Jacobian, slower-converging chord steps; tradeoff to evaluate.
  \item \textbf{Monotone projection is not the right path}. The Phi
        operator as defined doesn't admit a monotone FP. Either modify the
        operator definition (cusp-aware quadrature, or different lookup
        builder), or accept that the equilibrium is non-monotone.
\end{enumerate}

\section{Reproducibility}

\texttt{projects/REZN/solver\_code/highprec\_dd\_qd/dd\_k3\_cheby\_trf.py}
(TRF on Cheby-native), \texttt{dd\_k3\_method\_c.py} (Anderson + projection
for comparison).

\end{document}
